\documentclass[11pt]{article}
\usepackage[letterpaper,hmargin=.88in,vmargin=.82in]{geometry}
\usepackage{amsmath,amssymb,amsthm,mathtools}
\usepackage[T1]{fontenc}
\usepackage{lmodern,microtype}
\usepackage[hidelinks]{hyperref}
\hypersetup{pdftitle={One extra worker rules out the square-root welfare compromise}}
\newtheorem{theorem}{Theorem}
\newtheorem{lemma}{Lemma}
\newcommand{\E}{\mathbb E}
\newcommand{\Prob}{\mathbb P}
\newcommand{\ind}{\mathbf 1}
\newcommand{\eps}{\varepsilon_n}
\setlength{\parindent}{1.15em}
\setlength{\parskip}{2pt}
\raggedbottom
\begin{document}
\begin{center}
{\LARGE\bfseries One extra worker rules out the\\[2pt] square-root welfare compromise}
\end{center}

An unmatched worker can approach every firm. A firm that values her must
find a partner it values nearly as much. If every such partner earns low
utility at that firm, its assigned worker may in turn be willing to move
elsewhere. Can a clearinghouse still give both sides mean partner ranks
on the $\sqrt n$ scale, or both sides mean utility near $\tfrac12\log n$?

There are $n\ge2$ firms and $n+1$ workers. Firm $i$ values worker $j$ at $U_{ij}$,
and worker $j$ values firm $i$ at $V_{ij}$. All these values are independent
rate-one exponentials. An assignment matches each firm to a distinct worker.
A pair blocks if the firm would gain strictly more than
$\eps=\log n/(6\sqrt n)$ and the worker is either unmatched or would also
gain strictly more than $\eps$. An assignment is stable if no pair blocks.
Deferred acceptance supplies a zero-cost stable assignment, which is also
stable at this cost.

\begin{theorem}\label{thm:main}
Put $\delta=1/400000$. With probability tending to one over markets,
every stable assignment simultaneously satisfies
\[
  \overline R_W>n^{1/2+\delta},
  \qquad
  \overline V_W<(1/2-\delta)\log n .
\]
Here $\overline R_W$ and $\overline V_W$ are the mean partner rank and
mean utility among the $n$ matched workers. Rank one is best.
\end{theorem}

Draw the market first. Even a clearinghouse that sees all preferences and
chooses the stable assignment with the lowest worker mean rank misses the rank
benchmark. Choosing instead the one with the highest worker mean utility misses
the utility benchmark on the same event.

Over the finite stable set, let $C_n$ be the minimum of the larger of the two
sides' mean ranks among matched agents. Let $H_n$ be the maximum of the smaller
mean utility divided by $\log n$, counting the unmatched worker at zero among
all $n+1$ workers. Then $\log C_n/\log n>1/2+\delta$ and $H_n<1/2-\delta$ on
the theorem's event, ruling out both proposed limits of $1/2$.

\paragraph{From firm utility to both worker outcomes.}
Blocking threats force firms above the half-log utility benchmark
(Section~\ref{sec:firm}). To see what this leaves for workers, let
$R_{ij}$ be worker $j$'s rank for firm $i$ and put
$W_{ij}=\max\{V_{ij},\log(n/R_{ij})\}$. For every fixed $\eta>0$,
with probability tending to one, every assignment $M$ satisfies
\begin{equation}\label{eq:capacity}
  \frac1n\sum_i (U_{iM(i)}+W_{iM(i)}) < (1+\eta)\log n .
\end{equation}
Indeed, $R_{ij}$ is uniform on $\{1,\ldots,n\}$ by exchangeability and
absence of ties. In distribution $R_{ij}=\lceil nZ\rceil$ for $Z$ uniform
on $(0,1)$, so $\log(n/R_{ij})$ is stochastically bounded above by the
rate-one exponential $-\log Z$. For $0<\theta<1$, bounding the maximum
of the two exponentials by their sum gives
$\E e^{\theta W_{ij}}\le 2/(1-\theta)$. Independence of firm and worker
preferences gives $\E e^{\theta(U_{ij}+W_{ij})}\le 2/(1-\theta)^2$.
By Jensen's inequality on the assigned pairs, then enlarging the sum to
all pairs, a violation of \eqref{eq:capacity} by any $M$ implies
\[
 \frac1n\sum_{i,j}e^{\theta(U_{ij}+W_{ij})}\ \ge\ n^{\theta(1+\eta)}.
\]
The left side has mean at most $2(n+1)/(1-\theta)^2$. Taking
$1/(1+\eta)<\theta<1$ makes the probability of a violation tend to zero
by Markov's inequality.

Write $S(M)=\sum_i e^{-U_{iM(i)}}$. Section~\ref{sec:firm} proves
$S(M)\le n^{1/2-2\delta}$ uniformly over stable assignments with high
probability. Convexity of $-\log$ then gives mean firm utility at least
$\log(n/S(M))\ge(1/2+2\delta)\log n$. Apply \eqref{eq:capacity} with
$\eta=\delta$. The mean of $W$ on matched pairs is below
$(1/2-\delta)\log n$. It bounds both $\overline V_W$ and
$\log(n/\overline R_W)$ from above, the latter by convexity of $-\log$.
This proves Theorem~\ref{thm:main}. Including the unmatched worker multiplies
the nonnegative matched-worker mean utility by $n/(n+1)$.

\section{Why stability forces high firm utility}\label{sec:firm}

Use the tail probabilities $X_{ij}=e^{-U_{ij}}$ and $Y_{ij}=e^{-V_{ij}}$.
They are independent uniforms on $(0,1)$. Smaller values denote more
desirable partners. Write $q=e^{-\eps}$, so a strict gain of more than
$\eps$ means the new mark is less than $q$ times the current mark.
For an assignment let $x_i$ and $y_i$ be the two marks on firm $i$'s
assigned edge. Set $S=\sum_i x_i$, $T=\sum_i y_i$, $h=4\delta=10^{-5}$,
and $b_*=n^{-1/4}$.

\begin{lemma}[Workers willing to move]\label{lem:tail}
For all large $n$ and each $q\in[1/2,1]$, except on an event of probability
at most $(n+1)e^{-(\log n)^2/1024}$, every stable assignment satisfies
\[
    \sum_i y_i\,\ind_{\{y_i\ge b_*\}}\ \ge\ n^h-1 .
\]
\end{lemma}

A worker with mark $y_i$ would gain more than the switching cost at any
new firm with worker mark below $q y_i$. If firms also have a large total
mark $S$, these workers have many opportunities to form a blocking pair.
The appendix proves the worker bound. We now quantify that last step.

\begin{lemma}\label{lem:firm}
With probability tending to one, every stable assignment satisfies
$S\le n^{1/2-2\delta}$.
\end{lemma}

\begin{proof}
Set $s=n^{1/2-2\delta}$ and $B=n^h-1$. Count stable assignments with both
$S\ge s$ and worker tail sum at least $B$, and denote their number by $N$.
If some stable assignment has $S\ge s$, either some stable assignment has
tail sum below $B$ or $N\ge1$. Lemma~\ref{lem:tail} bounds the first
probability. The second is at most $\E N$.
Fix an injection and label its matched workers by their firms. Conditional
on its assigned marks, distinct off-assignment pairs use independent
edges. The probability of stability is
\[
 \prod_i(1-qx_i)\prod_{i\ne j}(1-q^2x_i y_j)
 \ \le\ \exp\!\left(-qS-q^2ST+q^2\sum_i x_i y_i\right)
 \ \le\ e^{-q^2ST}.
\]
The first product comes from the unmatched worker. Since $y_i\le1$ and
$q\le1$, it absorbs the correction for the assigned edges.

At fixed $S$, extend the worker integral to the positive orthant and put
$z_i=q^2 S y_i$. The factor $e^{-q^2ST}$ becomes the joint density of
independent rate-one exponentials $Z_i$ in the new variables, with
Jacobian $(q^2 S)^{-n}$. At a given firm sum $S$, write
$x_n=S-\sum_{i<n}x_i$. The other $n-1$ nonnegative coordinates
sum to at most $S$, a simplex of volume $S^{n-1}/(n-1)!$.
Thus the powers of $S$ leave just $dS/S$. There are $(n+1)!$ injections.
As $S$ increases from $s$, the required tail sum of the $Z_i$ increases
and fewer coordinates qualify. The expected count is therefore at most
\[
 n(n+1)q^{-2n}\log(n/s)\,
 \Prob\!\left(\sum_i Z_i\ind_{\{Z_i\ge q^2s b_*\}}\ge q^2sB\right).
\]
For $Z$ rate-one exponential and $a\ge0$, integration on $(0,a)$
and $[a,\infty)$ gives
$\E\exp(\tfrac12 Z\ind_{\{Z\ge a\}})=1-e^{-a}+2e^{-a/2}$.
Exponential Markov and independence bound the displayed probability by
$\exp[-q^2sB/2+2n e^{-q^2s b_*/2}]$. Now $s b_*=n^{1/4-2\delta}$
and $sB\sim n^{1/2+2\delta}$. Since $q\ge1/2$ eventually, the last
positive term tends to zero, and $2n\eps=\sqrt n\log n/3=o(sB)$.
Hence $\E N\to0$.
\end{proof}

\appendix
\section{How blocking threats spread}\label{sec:propagation}

To prove Lemma~\ref{lem:tail}, fix some assigned pairs and ask whether
\emph{any} stable completion can violate the bound.

Use scales $b_\ell=(\log n)^{-2}8192^{-\ell}$ for
$0\le\ell\le L=\lfloor\log n/(8\log8192)\rfloor$. Even the smallest
scale is at least $n^{-1/8}/(\log n)^2\ge b_*$ for large $n$.

Fix a candidate unmatched worker $u$ and an injective partial assignment $P$
excluding $u$. Fix both marks on every edge in the rows of the firms in $P$.
Let $J$ be its workers, $y_j$ their assigned worker marks, and put
$D=1+\sum_{j\in J}y_j$. The one represents the unmatched worker, who
is always willing to move. At scale $b=b_\ell$ assume that all recorded
$y_j\ge b$ and $D\ge(129/128)^\ell$. Write $E(P,u)$ for the event that some
stable completion of $P$ leaves $u$ unmatched and has tail sum $<n^h-1$.
For every fixed value of the recorded rows, we claim the remaining independent
rows give $\Prob(E(P,u))\le e^{-D/(1024b)}$.
If $D\ge n^h$, the recorded partners alone preclude such a completion.
This settles the last stage, since
$\log(129/128)/(8\log8192)>h$ (use $\log(1+1/128)\ge1/129$ and
$\log8192<10$). The floor in $L$ costs only a constant factor.

\paragraph{Which firms must take such a worker?}
Suppose $D<n^h$ and $\ell<L$. Put $t=1/(nb)$.
In each unused row $i$, consider two tests.
First, some old worker threatens the firm if it takes a partner with
$x_i>t$. This test passes if $X_{iu}<qt$, or if some $j\in J$ has
$X_{ij}<qt$ and $Y_{ij}<q y_j$. Second, no fresh worker
$j\notin J\cup\{u\}$ has both $X_{ij}\le t$ and $Y_{ij}\le b$.
Call the row forced if it passes both. In a stable completion, the first
test forces $x_i\le t$, since its witness is assigned elsewhere or
unmatched. The second then forces $Y_{iM(i)}>b$ for the actual partner $M(i)$.

The first test has probability
$1-(1-qt)\prod_{j\in J}(1-q^2t y_j)$ between $tD/8$ and $tD$.
For the lower bound use $qt+q^2t\sum_j y_j\ge tD/4$,
$1-z\le e^{-z}$, and $1-e^{-a}\ge a/2$ for $0\le a\le1$.
Here $tD\le n^{-3/4+h}$. The union bound gives the upper estimate.
The second test uses disjoint columns and has probability at least
$(1-1/n)^n\ge1/4$. There are at least $n/2$ unused rows because
$|J|\le D/b\le n^{h+1/4}$. The number of forced rows is binomial with
mean at least $D/(64b)$. The probability of fewer than $D/(128b)$ is
therefore at most $e^{-D/(512b)}$ by the half-mean Chernoff bound.

\paragraph{Paying for the choice of partners.}
If $E(P,u)$ occurs and there are at least $D/(128b)$ forced rows, take any
violating completion. Its partners on these rows form a partial assignment $F$
of new firms to distinct fresh workers. Each pair adds $y>b$, so
$D'=D+\sum_{F}y\ge(129/128)D$. For each fixed candidate $F$, retain the event
$A_F$ that its firms pass the first test, its specified edges have $X\le t$ and
$Y>b$, and $D'\ge(129/128)D$. Every violating completion on this branch
therefore witnesses $A_F\cap E(P\cup F,u)$ for at least one $F$.

For fixed $F$, reveal only its named rows. These determine $A_F$ given the
recorded rows. On $A_F$ the next stage's hypotheses hold. The skipped rows
retain their independent law because no test outcome on them was required.
Applying induction at scale $b/8192$ and averaging over the named rows gives
\[
 \Prob\big(A_F\cap E(P\cup F,u)\big)
 \ \le\ e^{-8D/b}\,\E\!\left[\ind_{A_F}\prod_{(i,j)\in F}e^{-8Y_{ij}/b}\right].
\]
On each chosen row the first test has probability at most $tD$ and uses
different columns from the specified partner's edge. That edge has $X\le t$
and $Y\in dy$ with probability $t\,dy$. After dropping the growth condition,
the row contributes at most $(tD)t\int_b^1e^{-8y/b}\,dy\le t^2Db/8$.
To bound the nonnegative sum over $F$, we can now also drop the requirement
that the chosen partners be distinct. Each unused row can be
skipped or can name one of at most $n$ fresh workers. Independence across
rows bounds the contribution of all valid choices by
\[
 e^{-8D/b}(1+nt^2Db/8)^n
 \ \le\ \exp[-(8-1/8)D/b],
 \qquad n^2t^2 b=1/b .
\]
Adding the probability of too few forced rows yields at most
$2e^{-D/(512b)}\le e^{-D/(1024b)}$, since $D/b\ge(\log n)^2$.
This closes the backward induction. Start with the
empty assignment, $D=1$, and sum over the $n+1$ possible identities of
$u$ to obtain Lemma~\ref{lem:tail}.\hfill$\square$
\end{document}
